\section{Randomized --- Biased Coin Halting}
%%% generated by the blueprint pipeline from TCSlib.Complexity.Randomized.BiasedCoinHalting
\subsection*{Randomized --- Biased Coin Halting: Introduction}

Declarations of the Lean module \texttt{TCSlib.Complexity.Randomized.BiasedCoinHalting}.

\subsection*{Randomized --- Biased Coin Halting: Results}

\begin{definition}[Halting index predicate]
\lean{Randomized.haltDigit}
\label{Randomized.haltDigit}
\leanok
\uses{Nat.Partrec.Code.haltingSet}
A natural number $j$ has the halting property if it belongs to the halting set of
natural-number partial recursive program codes.
\end{definition}

\begin{theorem}[Undecidability of the halting digit predicate]
\lean{Randomized.not_computablePred_haltDigit}
\label{Randomized.not_computablePred_haltDigit}
\leanok
\uses{Nat.Partrec.Code.not_computablePred_mem_haltingSet, Randomized.haltDigit}
\difficulty{1}
Let the halting set consist of codes of programs paired with inputs for which the
encoded program halts on the encoded input, and let the halting digit predicate hold of
exactly the natural numbers in this set. This predicate is not computable.
\end{theorem}

\begin{definition}[Halting indicator]
\lean{Randomized.haltBit}
\label{Randomized.haltBit}
\leanok
\uses{Randomized.haltDigit}
For each natural number $j$, define its halting indicator to be $1$ if $j$ belongs to
the halting set of natural-number partial recursive program codes, and $0$ otherwise.
\end{definition}

\begin{theorem}[Bound on the halting bit]
\lean{Randomized.haltBit_le_one}
\label{Randomized.haltBit_le_one}
\leanok
\uses{Randomized.haltBit}
\difficulty{3}
For every natural number $j$, the halting bit associated with $j$ is at most $1$.
\end{theorem}

\begin{definition}[Halting probability coin]
\lean{Randomized.haltCoin}
\label{Randomized.haltCoin}
\leanok
\uses{Randomized.haltBit}
Let $b_j$ be $1$ if the $j$th encoded computation halts and $0$ otherwise. The halting
coin is the real number $\sum_{j=0}^{\infty} b_j4^{-(j+1)}$.
\end{definition}

\begin{definition}[Halting digit prefixes]
\lean{Randomized.haltPrefix}
\label{Randomized.haltPrefix}
\leanok
\uses{Randomized.haltBit}
Let $b_j$ equal $1$ when the $j$th encoded computation halts and $0$ otherwise. Define
the natural number $P_J$ recursively by $P_0=0$ and $P_{J+1}=4P_J+b_J$.
\end{definition}

\begin{theorem}[Summability of the halting-bit series]
\lean{Randomized.summable_haltCoin}
\label{Randomized.summable_haltCoin}
\leanok
\uses{Randomized.haltBit, Randomized.haltBit_le_one}
\difficulty{4}
For every natural number $J$, the series over $j \in \mathbb{N}$ with terms
$h_{j+J}4^{-(j+1)}$ is summable, where $h_k=1$ if the $k$th natural number belongs to
the halting set and $h_k=0$ otherwise.
\end{theorem}

\begin{theorem}[Halting coin and finite prefixes]
\lean{Randomized.haltCoin_mul_pow}
\label{Randomized.haltCoin_mul_pow}
\leanok
\uses{Randomized.haltBit, Randomized.haltBit_le_one, Randomized.haltCoin, Randomized.haltPrefix, Randomized.summable_haltCoin}
\difficulty{6}
For every natural number $J$, there is a real number $\tau$ such that the halting coin
multiplied by $4^J$ equals the length-$J$ halting prefix plus $\tau$, with $0\le\tau\le
1/3$.
\end{theorem}

\begin{theorem}[Bounds on the halting coin]
\lean{Randomized.haltCoin_mem}
\label{Randomized.haltCoin_mem}
\leanok
\uses{Randomized.haltCoin, Randomized.haltCoin_mul_pow, Randomized.haltPrefix}
\difficulty{3}
The halting coin lies in the interval $[0,1/3]$.
\end{theorem}

\begin{definition}[Length of the binary arithmetic head]
\lean{Randomized.hlN}
\label{Randomized.hlN}
\leanok
\uses{Complexity.BinArith.uHead}
For a finite Boolean list $x$, write $h(x)$ for its head component under the
binary-arithmetic decomposition. Define the associated natural number to be the length
of $h(x)$.
\end{definition}

\begin{definition}[Base-four remainder-length parameter]
\lean{Randomized.hlJ}
\label{Randomized.hlJ}
\leanok
\uses{Complexity.BinArith.uRest}
For a finite Boolean list $x$, let $r(x)$ denote the remainder list in its
binary-arithmetic decomposition. Define $h_J(x)$ to be the base-$4$ logarithm of the
length of $r(x)$, rounded down to an integer, with value $0$ when that length is $0$.
\end{definition}

\begin{definition}[Halting language]
\lean{Randomized.haltLang}
\label{Randomized.haltLang}
\leanok
\uses{Randomized.haltDigit, Randomized.haltPrefix, Randomized.hlJ, Randomized.hlN}
For a finite Boolean word $x$, let $J(x)$ and $N(x)$ be the natural-number quantities
assigned to it, and let $P(k)$ and $D(k)$ be the associated prefix and digit predicates.
The language consists of exactly those $x$ for which $J(x)\ge 1$, $N(x)=P(J(x)-1)$, and
$D(J(x)-1)$ holds.
\end{definition}

\begin{theorem}[Bounds for the halting-number approximation]
\lean{Randomized.haltCoin_mul_pow_hlJ}
\label{Randomized.haltCoin_mul_pow_hlJ}
\leanok
\uses{Randomized.haltBit, Randomized.haltBit_le_one, Randomized.haltCoin, Randomized.haltCoin_mul_pow, Randomized.haltDigit, Randomized.haltLang, Randomized.haltPrefix, Randomized.hlJ, Randomized.hlN}
\difficulty{6}
For every finite Boolean string $x$, let $N$ be the length of the output of the head
transducer on $x$, and let $J$ be the base-$4$ logarithm, rounded down, of the length of
the output of the rest transducer on $x$. If $x$ belongs to the language of strings
satisfying $J \ge 1$, $N$ equal to the halting-bit prefix through position $J-1$, and a
halting bit at position $J-1$, then $4N+1 \le H4^J \le 4N+4/3$, where $H$ is the halting
coin. If $x$ does not belong to that language, then $H4^J \le 4N+1/3$ or $4N+4 \le
H4^J$.
\end{theorem}

\begin{definition}[Halting query string]
\lean{Randomized.hlQuery}
\label{Randomized.hlQuery}
\leanok
For natural numbers $N$ and $j$, let $q$ be the Boolean list formed by concatenating $N$
true values, one false value, and $4^{j+1}$ true values, in that order.
\end{definition}

\begin{theorem}[Parameters of the query string]
\lean{Randomized.hlN_hlQuery}
\label{Randomized.hlN_hlQuery}
\leanok
\uses{BoolCircuit.encodeNat, Complexity.BinArith.uHead_encodeNat, Complexity.BinArith.uRest_encodeNat, Randomized.hlJ, Randomized.hlN, Randomized.hlQuery}
\difficulty{4}
For all natural numbers $N$ and $j$, the query string consisting of $N$ ones, followed
by a zero and then $4^{j+1}$ ones, has head-output length $N$ and rest-output
logarithmic parameter $j+1$.
\end{theorem}

\begin{theorem}[Query Membership and Halting]
\lean{Randomized.hlQuery_mem_iff}
\label{Randomized.hlQuery_mem_iff}
\leanok
\uses{Randomized.haltDigit, Randomized.haltLang, Randomized.haltPrefix, Randomized.hlJ, Randomized.hlN, Randomized.hlN_hlQuery, Randomized.hlQuery}
\difficulty{3}
For every natural number $j$, the word consisting of $\operatorname{haltPrefix}(j)$
ones, followed by a zero and then $4^{j+1}$ ones, belongs to the language exactly when
the computation indexed by $j$ halts.
\end{theorem}

\begin{theorem}[Undecidability of the halting language]
\lean{Randomized.not_computablePred_haltLang}
\label{Randomized.not_computablePred_haltLang}
\leanok
\uses{Randomized.haltBit, Randomized.haltDigit, Randomized.haltLang, Randomized.haltPrefix, Randomized.hlQuery, Randomized.hlQuery_mem_iff, Randomized.not_computablePred_haltDigit}
\difficulty{6}
Let the halting set consist of codes of programs paired with inputs for which the
encoded program halts on the encoded input. Define the halting digit to be $1$ for
numbers in this set and $0$ otherwise, and define the halting prefix by $h(0)=0$ and
$h(J+1)=4h(J)+	ext{halting digit}(J)$. For a binary string $x$, let $N(x)$ be the length
of the output of the specified head transducer on $x$, and let $J(x)$ be the base-$4$
logarithm, rounded down, of the length of the output of the specified rest transducer on
$x$. The language consists of strings $x$ satisfying $1\le J(x)$, $N(x)=h(J(x)-1)$, and
membership of $J(x)-1$ in the halting set; this language is not computable.
\end{theorem}

\begin{definition}[Fourth power of a list parameter]
\lean{Randomized.hlD}
\label{Randomized.hlD}
\leanok
\uses{Randomized.hlJ}
For each finite Boolean list $x$, let $h(x)$ denote the natural-number parameter
assigned to $x$. Define $D(x)=4^{h(x)}$.
\end{definition}

\begin{definition}[Halt coin verifier]
\lean{Randomized.haltCoinVerifier}
\label{Randomized.haltCoinVerifier}
\leanok
\uses{Randomized.hlD, Randomized.hlN}
For Boolean lists $x$ and $r$, the verifier returns true exactly when both inequalities
hold: $|r|(12N+2)<3T D$ and $T D<|r|(4N+2)$, where $N$ and $D$ are the numerator and
denominator quantities associated with the halting problem for $x$, and $T$ is the
number of true entries in $r$.
\end{definition}

\begin{theorem}[Coin verifier identifies the halting language]
\lean{Randomized.haltCoinVerifier_eq}
\label{Randomized.haltCoinVerifier_eq}
\leanok
\uses{Randomized.haltCoin, Randomized.haltCoinVerifier, Randomized.haltCoin_mul_pow_hlJ, Randomized.haltLang, Randomized.hlD, Randomized.hlJ, Randomized.hlN, Randomized.indicator_eq_false_iff, Randomized.indicator_eq_true_iff, Turing.MultiTapeTM.indicator}
\difficulty{6}
Let $x$ and $r$ be finite bit strings. Suppose $r$ has positive length and the fraction
of its entries equal to $1$ differs from the halting coin by at most $1/(4d(x))$, where
$d(x)$ is the power of $4$ determined by the length parameter of $x$. Then the
verifier's output on $(x,r)$ equals the Boolean indicator of whether $x$ belongs to the
halting language.
\end{theorem}

\begin{theorem}[Bound on the depth parameter]
\lean{Randomized.hlD_le}
\label{Randomized.hlD_le}
\leanok
\uses{Complexity.BinArith.length_uRest_le, Complexity.BinArith.uRest, Randomized.hlD, Randomized.hlJ}
\difficulty{4}
For every finite binary string $x$, let $u$ be the output of the specified one-pass
transducer on $x$, and set $j=\lfloor\log_4 |u|\rfloor$ (with the natural-number
logarithm convention). Then $4^j\le |x|+1$.
\end{theorem}

\begin{theorem}[Success probability of the halting-coin verifier]
\lean{Randomized.haltCoinVerifier_prob}
\label{Randomized.haltCoinVerifier_prob}
\leanok
\uses{Randomized.biasedProb, Randomized.biasedProb_count_dev_le, Randomized.biasedProb_mono, Randomized.biasedProb_not, Randomized.haltCoin, Randomized.haltCoinVerifier, Randomized.haltCoinVerifier_eq, Randomized.haltCoin_mem, Randomized.haltLang, Randomized.hlD, Randomized.hlD_le, Randomized.polyLen, Turing.MultiTapeTM.indicator}
\difficulty{6}
For every finite bit string $x$, choose a random bit string of length $32(|x|+1)^2$ by
independently setting each bit to $1$ with probability equal to the halting coin and to
$0$ otherwise. The probability that the verifier's output equals the Boolean indicator
of membership of $x$ in the halting language is at least $2/3$.
\end{theorem}

\begin{definition}[Comparison of paired list lengths]
\lean{Randomized.lenLeB}
\label{Randomized.lenLeB}
\leanok
\uses{Complexity.pairFstD, Complexity.pairSndD}
For a Boolean list $z$, return true exactly when the length of its second decoded
component is at most the length of its first decoded component.
\end{definition}

\begin{definition}[Conditional fourth-power step]
\lean{Randomized.pow4Step}
\label{Randomized.pow4Step}
\leanok
\uses{Complexity.pairFstD, Complexity.pairSndD, Randomized.lenLeB, Turing.pairEncode}
For a Boolean list $z$, encode its first decoded component together with a second
component chosen as follows: use four concatenated copies of its second decoded
component if, after encoding the first component with those four copies, the second
decoded component has length at most the first; otherwise use the original second
decoded component.
\end{definition}

\begin{theorem}[Length growth under iteration]
\lean{Randomized.iterate_pow4Step}
\label{Randomized.iterate_pow4Step}
\leanok
\uses{Complexity.pairFstD_pairEncode, Complexity.pairSndD_pairEncode, Randomized.lenLeB, Randomized.pow4Step, Turing.pairEncode}
\difficulty{6}
For every bit string $m$ and natural number $t$, there is a bit string $w$ such that
after applying the fourth-power step $t$ times to the encoding of $m$ paired with
$[1]$, the result is the encoding of $m$ paired with $w$, and
$|w|=4^{\min(t,\operatorname{Nat.log}(4,|m|))}$.
\end{theorem}

\begin{definition}[Halting-word transformation]
\lean{Randomized.hlDWord}
\label{Randomized.hlDWord}
\leanok
\uses{Complexity.BinArith.forLoop, Complexity.BinArith.uRest, Complexity.pairSndD, Randomized.pow4Step, Turing.pairEncode}
For a Boolean list $x$, run the loop defined by the fourth-power step on the encoding of
the pair whose first component is the remainder of $x$ after removing its unary prefix
and whose second component encodes that same remainder together with the one-element
list containing true. Then take the second decoded component twice successively.
\end{definition}

\begin{theorem}[Length of the iterated output word]
\lean{Randomized.length_hlDWord}
\label{Randomized.length_hlDWord}
\leanok
\uses{Complexity.BinArith.forLoop_pairEncode, Complexity.BinArith.uRest, Complexity.pairSndD_pairEncode, Randomized.hlD, Randomized.hlDWord, Randomized.hlJ, Randomized.iterate_pow4Step}
\difficulty{3}
For every bit string $x$, let $u$ be the word obtained by applying the specified
one-pass binary transducer to $x$, and let $W(x)$ be the word obtained from $u$ by the
specified pairing, iteration, and projection construction. Then $|W(x)|=4^{\lfloor\log_4
|u|\rfloor}$.
\end{theorem}

\begin{definition}[Positions of true bits]
\lean{Randomized.onesW}
\label{Randomized.onesW}
\leanok
\uses{Complexity.transduce}
For a Boolean list $w$, return the list produced by transducing $w$ with a one-state
machine that emits true for each true input bit and emits nothing for each false input
bit.
\end{definition}

\begin{theorem}[Length of the true-bit output]
\lean{Randomized.length_onesW}
\label{Randomized.length_onesW}
\leanok
\uses{Complexity.transduce, Randomized.onesW}
\difficulty{4}
For every Boolean word $w$, the word formed by deleting its false bits has length equal
to the number of true bits in $w$.
\end{theorem}

\begin{definition}[Linear-size encoding]
\lean{Randomized.hcLin}
\label{Randomized.hcLin}
\leanok
\uses{Complexity.BinArith.mulU, Complexity.BinArith.uHead, Complexity.pairFstD, Complexity.pairSndD, Turing.pairEncode}
For a natural number $n$ and a Boolean list $z$, encode the pair consisting of the
second decoded component of $z$ and the concatenation of two true bits with the result
of multiplying, in unary, the encoding of the list of $n$ true bits and the unary head
of the first decoded component of $z$. Then multiply this encoded list in unary.
\end{definition}

\begin{definition}[Conditional-description encoding]
\lean{Randomized.hcCD}
\label{Randomized.hcCD}
\leanok
\uses{Complexity.BinArith.mulU, Complexity.pairFstD, Complexity.pairSndD, Randomized.hlDWord, Randomized.onesW, Turing.pairEncode}
For a Boolean list $z$, take the unary product of the encoding of the pair consisting of
the true-bit positions in its second decoded component and the word obtained by applying
the halting-word transformation to its first decoded component.
\end{definition}

\begin{definition}[Three-bit conditional-description encoding]
\lean{Randomized.hcCD3}
\label{Randomized.hcCD3}
\leanok
\uses{Complexity.BinArith.mulU, Complexity.pairFstD, Complexity.pairSndD, Randomized.hlDWord, Randomized.onesW, Turing.pairEncode}
For a Boolean list $z$, take the unary product of the encoding of the pair consisting of
the true-bit positions in its second decoded component and the unary product of the
encoding of the pair formed by the halting-word transformation of its first decoded
component and the three-element list containing only true bits.
\end{definition}

\begin{definition}[Halting coin test]
\lean{Randomized.haltCoinTest}
\label{Randomized.haltCoinTest}
\leanok
\uses{Randomized.hcCD, Randomized.hcCD3, Randomized.hcLin, Randomized.lenLeB, Turing.pairEncode}
For a Boolean list $z$, return true exactly when both of the following hold: the second
decoded component of the encoding of the linear-size encoding of $z$ at parameter $12$
and its three-bit conditional-description encoding has length greater than the first
decoded component's length; and the second decoded component of the encoding of its
conditional-description encoding and the linear-size encoding of $z$ at parameter $4$
has length greater than the first decoded component's length.
\end{definition}

\begin{theorem}[Pair-encoded coin test identity]
\lean{Randomized.haltCoinTest_pairEncode}
\label{Randomized.haltCoinTest_pairEncode}
\leanok
\uses{Complexity.BinArith.length_mulU, Complexity.pairFstD_pairEncode, Complexity.pairSndD_pairEncode, Randomized.haltCoinTest, Randomized.haltCoinVerifier, Randomized.hcCD, Randomized.hcCD3, Randomized.hcLin, Randomized.hlN, Randomized.lenLeB, Randomized.length_hlDWord, Randomized.length_onesW, Turing.pairEncode}
\difficulty{1}
For any two finite bit strings $x$ and $r$, applying the halting-coin test to their pair
encoding gives the same Boolean value as the verifier applied to $x$ and $r$.
\end{theorem}

\begin{theorem}[Polynomial-time fourth-power string map]
\lean{Randomized.polyTimeComputable_pow4Step}
\label{Randomized.polyTimeComputable_pow4Step}
\leanok
\uses{Complexity.PolyTimeComputable, Complexity.PolyTimeComputable.append, Complexity.PolyTimeComputable.comp, Complexity.PolyTimeComputable.pairEncode, Complexity.pairFstD, Complexity.pairSndD, Complexity.polyTimeComputable_ite, Complexity.polyTimeComputable_lenLe, Complexity.polyTimeComputable_pairFstD, Complexity.polyTimeComputable_pairSndD, Randomized.lenLeB, Randomized.pow4Step, Turing.pairEncode}
\difficulty{4}
For a bit string $z$, let $a$ and $b$ be its decoded first and second components, taking
either to be empty when decoding fails. Form the pair whose first component is $a$ and
whose second component is four concatenated copies of $b$ if the length of that repeated
string is at most the length of $a$; otherwise use $b$ as the second component. This
function is polynomial-time computable.
\end{theorem}

\begin{theorem}[Polynomial-Time Length-Limited Word Construction]
\lean{Randomized.polyTimeComputable_hlDWord}
\label{Randomized.polyTimeComputable_hlDWord}
\leanok
\uses{Complexity.BinArith.forLoop_bound_of_ctx, Complexity.BinArith.polyTimeComputable_forLoop, Complexity.BinArith.polyTimeComputable_slot3, Complexity.BinArith.polyTimeComputable_uRest, Complexity.BinArith.uRest, Complexity.PolyTimeComputable, Complexity.PolyTimeComputable.comp, Complexity.PolyTimeComputable.pairEncode, Complexity.length_pair_components_le, Complexity.pairFstD, Complexity.pairFstD_pairEncode, Complexity.pairSndD, Complexity.pairSndD_pairEncode, Complexity.polyTimeComputable_const, Complexity.polyTimeComputable_pairFstD, Randomized.hlDWord, Randomized.lenLeB, Randomized.polyTimeComputable_pow4Step, Randomized.pow4Step, Turing.length_pairEncode}
\difficulty{5}
For every bit string $z$, the function that applies the length-limited fourfold
iteration to the tail transducer output on the first component of the decoded pair is
computable by a deterministic finite-state Turing machine in time $C(n+1)^c$ on every
input of length $n$, for some natural numbers $C$ and $c$.
\end{theorem}

\begin{theorem}[Polynomial-Time Parameterized Linear Construction]
\lean{Randomized.polyTimeComputable_hcLin}
\label{Randomized.polyTimeComputable_hcLin}
\leanok
\uses{Complexity.BinArith.mulU, Complexity.BinArith.polyTimeComputable_app, Complexity.BinArith.polyTimeComputable_mulU, Complexity.BinArith.polyTimeComputable_uHead, Complexity.BinArith.uHead, Complexity.PolyTimeComputable, Complexity.PolyTimeComputable.append, Complexity.PolyTimeComputable.pairEncode, Complexity.pairFstD, Complexity.pairSndD, Complexity.polyTimeComputable_const, Complexity.polyTimeComputable_pairFstD, Complexity.polyTimeComputable_pairSndD, Randomized.hcLin, Turing.pairEncode}
\difficulty{4}
For every natural number $n$, the function that, on a bit string $z$, decodes its pair
and applies the unary multiplication construction twice—first to the pair consisting of
the word of $n$ ones and the head transducer output on the first component, then to the
pair consisting of the decoded second component and the first multiplication output
followed by $11$—is computable by a deterministic finite-state Turing machine in time
$C(m+1)^c$ on every input of length $m$, for some natural numbers $C$ and $c$.
\end{theorem}

\begin{theorem}[Polynomial-time second-component filtering]
\lean{Randomized.polyTimeComputable_onesW_snd}
\label{Randomized.polyTimeComputable_onesW_snd}
\leanok
\uses{Complexity.BinArith.polyTimeComputable_app, Complexity.PolyTimeComputable, Complexity.pairSndD, Complexity.polyTimeComputable_pairSndD, Complexity.polyTimeComputable_transduce, Randomized.onesW}
\difficulty{1}
For a bit string $z$, decode its second component using the aligned pair parser,
returning the empty string if decoding fails, and then delete every $0$ from that
component. The resulting function of $z$ is polynomial-time computable.
\end{theorem}

\begin{theorem}[Polynomial-Time Coin Construction]
\lean{Randomized.polyTimeComputable_hcCD}
\label{Randomized.polyTimeComputable_hcCD}
\leanok
\uses{Complexity.BinArith.polyTimeComputable_app, Complexity.BinArith.polyTimeComputable_mulU, Complexity.PolyTimeComputable, Complexity.PolyTimeComputable.pairEncode, Complexity.pairFstD, Complexity.pairSndD, Randomized.hcCD, Randomized.hlDWord, Randomized.onesW, Randomized.polyTimeComputable_hlDWord, Randomized.polyTimeComputable_onesW_snd, Turing.pairEncode}
\difficulty{3}
For every bit string $z$, decode its pair, retain the $1$ bits of the second component
in their original order, and combine this string with the word obtained by applying the
fourfold length-limited iteration to the one-pass transducer output on the first
component; then apply the unary multiplication construction to the resulting pair. This
function is computable by a deterministic finite-state Turing machine in time $C(n+1)^c$
on every input of length $n$, for some natural numbers $C$ and $c$.
\end{theorem}

\begin{theorem}[Polynomial-Time Threefold Coin Construction]
\lean{Randomized.polyTimeComputable_hcCD3}
\label{Randomized.polyTimeComputable_hcCD3}
\leanok
\uses{Complexity.BinArith.mulU, Complexity.BinArith.polyTimeComputable_app, Complexity.BinArith.polyTimeComputable_mulU, Complexity.PolyTimeComputable, Complexity.PolyTimeComputable.pairEncode, Complexity.pairFstD, Complexity.pairSndD, Complexity.polyTimeComputable_const, Randomized.hcCD3, Randomized.hlDWord, Randomized.onesW, Randomized.polyTimeComputable_hlDWord, Randomized.polyTimeComputable_onesW_snd, Turing.pairEncode}
\difficulty{4}
For every bit string $z$, decode its pair, retain the $1$ bits of the second component
in their original order, and combine this string with the unary multiplication
construction applied to the pair consisting of the fourfold length-limited iteration of
the one-pass transducer output on the first component and the word $111$; then apply the
unary multiplication construction to the resulting pair. This function is computable by
a deterministic finite-state Turing machine in time $C(n+1)^c$ on every input of length
$n$, for some natural numbers $C$ and $c$.
\end{theorem}

\begin{theorem}[Polynomial-Time Halt-Coin Test]
\lean{Randomized.polyTimeComputable_haltCoinTest}
\label{Randomized.polyTimeComputable_haltCoinTest}
\leanok
\uses{Complexity.BinArith.polyTimeComputable_test_app, Complexity.PolyTimeComputable, Complexity.PolyTimeComputable.pairEncode, Complexity.polyTimeComputable_and, Complexity.polyTimeComputable_lenLe, Complexity.polyTimeComputable_not, Randomized.haltCoinTest, Randomized.hcCD, Randomized.hcCD3, Randomized.hcLin, Randomized.lenLeB, Randomized.polyTimeComputable_hcCD, Randomized.polyTimeComputable_hcCD3, Randomized.polyTimeComputable_hcLin, Turing.pairEncode}
\difficulty{4}
For every bit string $z$, evaluate the Boolean test that is true exactly when both of
the following hold: the length of the second component of the encoded pair formed from
the output of the twelve-fold linear construction and the threefold coin construction
exceeds the length of its first component; and the length of the second component of the
pair formed from the coin construction and the fourfold linear construction does not
exceed the length of its first component. The function sending $z$ to the singleton list
containing this test result is computable by a deterministic finite-state Turing machine
in time $C(n+1)^c$ on every input of length $n$, for some natural numbers $C$ and $c$.
\end{theorem}

\begin{theorem}[Efficiency of the halting coin verifier]
\lean{Randomized.eff_haltCoinVerifier}
\label{Randomized.eff_haltCoinVerifier}
\leanok
\uses{Randomized.boolVerifier, Randomized.eff_boolVerifier_of_test, Randomized.haltCoinTest_pairEncode, Randomized.haltCoinVerifier, Randomized.polyTimeComputable_haltCoinTest, Randomized.polyTimeModel}
\difficulty{1}
The verifier that accepts exactly when the fraction of true bits in its random string
lies strictly between the two bounds determined by the input’s head length and depth
runs in polynomial time in the verifier model.
\end{theorem}

\begin{theorem}[Halting language in biased randomized polynomial time]
\lean{Randomized.haltLang_mem_biasedBPP}
\label{Randomized.haltLang_mem_biasedBPP}
\leanok
\uses{Randomized.biasedBPP, Randomized.eff_haltCoinVerifier, Randomized.haltCoin, Randomized.haltCoinVerifier_prob, Randomized.haltLang, Randomized.mem_biasedBPP_of_eff}
\difficulty{1}
The language of Boolean strings whose logarithmic length parameter is at least $1$,
whose output length equals the corresponding prefix of the halting bits, and whose
preceding halting bit is $1$ belongs to the class of languages decidable in polynomial
time by a finite-state randomized Turing machine with independent choice bits that are
true with probability equal to the halting number, and with acceptance probability at
least $2/3$ on every input.
\end{theorem}

\begin{theorem}[Biased randomized computation beyond computability]
\lean{Randomized.exists_biasedBPP_not_computablePred}
\label{Randomized.exists_biasedBPP_not_computablePred}
\leanok
\uses{Randomized.biasedBPP, Randomized.haltCoin, Randomized.haltCoin_mem, Randomized.haltLang, Randomized.haltLang_mem_biasedBPP, Randomized.not_computablePred_haltLang}
\difficulty{3}
There exist a real number $\rho$ with $0\le \rho\le 1$ and a language $L$ of bit strings
that is not computable, yet belongs to the bounded-error polynomial-time class using
independent random bits that equal $1$ with probability $\rho$.
\end{theorem}
